% !TEX root = ../main.tex
% ============================================================================
%                    VARIABEL METADATA TESIS
% ============================================================================
% PETUNJUK: Ganti placeholder {} dengan data Anda
%
% Contoh pengisian:
% \newcommand{\ThesisTitle}{JUDUL TESIS ANDA DENGAN HURUF KAPITAL}
% \newcommand{\AuthorName}{NAMA LENGKAP ANDA}
% ============================================================================

% --- Data Judul dan Penulis ---
% {Judul Tesis Lengkap}: Judul tesis dalam huruf kapital
\newcommand{\ThesisTitle}{{Judul Tesis Lengkap}}

% {NAMA LENGKAP PENULIS}: Nama lengkap dalam huruf kapital
\newcommand{\AuthorName}{{NAMA LENGKAP PENULIS}}

% {Nama Lengkap Penulis}: Nama lengkap dengan huruf kapital di awal kata
\newcommand{\AuthorNameB}{{Nama Lengkap Penulis}}

% {NIM}: Nomor Induk Mahasiswa, contoh: P31.2024.02610
\newcommand{\StudentID}{{NIM}}

% --- Data Program Studi (Tidak perlu diubah) ---
\newcommand{\ProgramName}{PROGRAM PASCASARJANA\\MAGISTER TEKNIK INFORMATIKA\\UNIVERSITAS DIAN NUSWANTORO}

% --- Data Lokasi dan Waktu ---
% {Kota}: Kota tempat pembuatan tesis
\newcommand{\City}{{Kota}}

% {Tahun}: Tahun pembuatan tesis
\newcommand{\Year}{{Tahun}}

% {Tujuan Tesis}: Kalimat tujuan tesis
\newcommand{\ThesisPurpose}{Tesis diajukan sebagai salah satu syarat untuk memperoleh gelar Magister Komputer}

% {Alamat Penulis}: Alamat tetap penulis
\newcommand{\AuthorAddress}{{Alamat Penulis}}

% --- Data Pembimbing ---
% {Nama Pembimbing Utama}: Nama lengkap dengan gelar
% Contoh: Prof. Dr. Nama Lengkap, S.T., M.Kom.
\newcommand{\SupervisorOne}{{Nama Pembimbing Utama}}

% {NIDN Pembimbing Utama}: Nomor Induk Dosen Nasional
\newcommand{\SupervisorOneNIDN}{{NIDN Pembimbing Utama}}

% {Nama Pembimbing Kedua}: Nama lengkap dengan gelar
\newcommand{\SupervisorTwo}{{Nama Pembimbing Kedua}}

% {NIDN Pembimbing Kedua}: Nomor Induk Dosen Nasional
\newcommand{\SupervisorTwoNIDN}{{NIDN Pembimbing Kedua}}

% --- Data Dewan Penguji ---
% {Nama Penguji 1}: Nama lengkap dengan gelar
\newcommand{\ExaminerOne}{{Nama Penguji 1}}

% {Nama Penguji 2}: Nama lengkap dengan gelar
\newcommand{\ExaminerTwo}{{Nama Penguji 2}}

% {Nama Dekan}: Nama lengkap dengan gelar
\newcommand{\DeanName}{{Nama Dekan}}

% {Nama Ketua Penguji}: Nama lengkap dengan gelar
\newcommand{\ChairExaminer}{{Nama Ketua Penguji}}
